\section{Introduction}\label{sec:introduction}

A stationary integral varifold is a measure-theoretic minimal surface with
integer multiplicity. Its tangent plane exists almost everywhere, but a tangent
plane alone need not describe an entire neighborhood. The basic regularity
question is whether the points without a smooth minimal neighborhood can carry
positive measure. We prove that they have zero $m$-dimensional measure.

Let $m,n\geq1$, let $U\subset\R^{m+n}$ be open, and let $E\subset U$ be
countably $m$-rectifiable. Let $\theta:E\to\N\setminus\{0\}$ be measurable and
locally integrable with respect to $\HH^m$. Write
\[
 \mu(A)=\int_{E\cap A}\theta\,\dd\HH^m,
 \qquad P(x)=P_{T_xE}.
\]
The associated integral varifold $V$ is stationary if
\begin{equation}\label{eq:stationarity}
 \int \tr(P(x)DX(x))\,\dd\mu(x)=0
 \quad\text{for every }X\in C_c^1(U;\R^{m+n}).
\end{equation}
A point $p\in\spt\mu\cap U$ is \emph{regular} if there are a neighborhood
$W$ of $p$, a smooth embedded minimal $m$-submanifold $\Sigma\subset W$, and
an integer $q\geq1$ such that the restriction of the varifold to $W$ is
$q|\Sigma|$. The remaining support points form $\operatorname{Sing}V$.

\begin{theorem}\label{thm:main}
Every stationary integral $m$-varifold in an open subset of $\R^{m+n}$ satisfies
\[
 \HH^m(\operatorname{Sing}V)=0.
\]
Here $m,n\geq1$ are arbitrary. No stability, minimizing, orientability,
density upper bound, or codimension-one hypothesis is required.
\end{theorem}

\begin{corollary}[Round spheres]\label{cor:round-spheres}
Let $m,n\geq1$ and let $V$ be an integral $m$-varifold on the full unit round
sphere $S^{m+n}$, stationary for variations tangent to the sphere. A support
point is regular if on a spherical neighborhood $V=q|\Sigma|$ for an integer
$q\geq1$ and a smooth embedded minimal $m$-submanifold $\Sigma$ of the sphere.
If $\operatorname{Sing}_S V$ denotes the remaining support points, then
\[
 \HH^m_{g_{\mathrm{round}}}(\operatorname{Sing}_S V)=0.
\]
No fixed upper bound on mass, density, or multiplicity is required, nor are
stability, orientability, or mod-two cyclicity. The same conclusion holds on
every round sphere of positive radius.
\end{corollary}

\Cref{sec:spheres} derives this corollary from \Cref{thm:main} by a
separate cone argument. That proof adapts the hypersurface first-variation
calculation in \cite[Lemma~6.1]{OpenAISolomonYau2026} and rederives it in
the present dimensions.

\subsection{Background and significance}

Allard's regularity theory gives a relatively open dense regular set and, in
particular, regularity at density-one points \cite{Allard1972}; see also
\cite[Theorem~23]{Allard1987} and
\cite[Chapter~5, Theorem~6.1 and Corollary~6.3]{Simon2018}. Higher integer
density presents a different issue: dividing a varifold by its density at a
point need not preserve the almost-everywhere lower density bound required by
that theorem in a neighborhood. Standard accounts of the first variation,
monotonicity, and integral varifolds may be found in \cite{Simon1983}.

Simon's 2018 notes already record the question whether the singular set
must be $\HH^m$-null even when the generalized mean curvature is zero
\cite[Chapter~5, Remark~6.4]{Simon2018}. The measure-zero singular-set
assertion is Conjecture~1.2 of Brena, Decio,
and De Lellis~\cite{BDD2025}. \Cref{thm:main} proves this conjecture in its
full Euclidean scope. They isolate a local criterion in Conjecture~1.3:
a support point with a unique tangent varifold $Q|\Pi|$, where $\Pi$ is an
$m$-plane and $Q$ a positive integer, should be regular when the mass of
$\{\theta\ne Q\}$ in a ball of radius $r$ about it is $o(r^m)$.
They then propose approximating the varifold by a smooth minimal graph
with the squared-height integral vanishing to every order, followed by a
unique-continuation step that upgrades this contact to local support containment
\cite[Conjectures~1.3--1.5]{BDD2025}. Their analysis of multiplicity and
height in flat cylinders also provides useful precedents for the local
estimates developed here \cite[Theorem~1.9 and Lemma~2.3]{BDD2025}.

Related progress concerns the differentiability of a countable covering.
Menne proved second-order rectifiability for integral varifolds of
locally bounded first variation \cite[Theorem~3.6]{Menne2013}. For stationary
integral varifolds of dimension at least two, Brena, De Lellis, and
Franceschini proved smooth rectifiability and infinite-order approximation
at $\mu$-almost every point by one smooth graph
\cite[Theorems~1.2 and~1.3]{BDLF2025}. The approximating graph in the latter
statement is not asserted to be minimal; Brena--Decio--De Lellis present
this as a weaker form of their approximation conjecture
\cite[Theorem~1.8]{BDD2025}. Kolasi\'nski obtains rectifiability of every
finite order and an alternative proof of the smooth-rectifiability
conclusion, following an architecture he credits to Menne
\cite[Theorem~1.1, Corollary~10.8, and the historical note in Section~1]
{Kolasinski2026}.
Smooth rectifiability gives a countable smooth covering outside a
$\mu$-null set; \Cref{thm:main} requires equality with one smooth minimal
submanifold throughout a neighborhood of almost every point.

In codimension one, Hirsch and Spolaor prove almost-everywhere regularity
for integral currents that are stationary as varifolds and have density
at most two at every support point \cite[Theorem~1.1]{HS2025}.
These current and multiplicity hypotheses distinguish that result from
the arbitrary-codimension statement here.

De Lellis, Hirsch, Lihn, and Spolaor construct, in a smooth Riemannian
metric, a stationary varifold with a double-plane tangent at an exceptional
point and topology accumulating there \cite[Theorem~1.1]{DLHLS2026}.\footnote{In
their respective disclosures, Kolasi\'nski reports using Claude AI and
Claude Code, mostly Opus~5; the abstract of De Lellis--Hirsch--Lihn--Spolaor
credits a crucial suggestion to \texttt{gpt-astra-6}. These attributions
follow the cited papers' own descriptions.}

\subsection{The argument}

The main estimate compares projected mass and tilt with their leading terms
kept signed. In product coordinates $\R^m\times\R^n$, let $\pi$ be
projection onto the first factor, fix a positive integer $Q$, and let
$P_{vv}$ be the vertical block of the tangent projection. With lengths
and mass normalized to the working radius, define
\[
 M=\pi_\#\mu-Q\,\dd y,
 \qquad D=\pi_\#(\tr P_{vv}\,\mu).
\]
Here $|M|$ denotes the total variation of the signed measure $M$, and
$\tau=|M|+D$ is the positive excess measure in this flat setting.
When $\tau$ has mass at most $d$ on a fixed enlarged ball, we prove an
interior bound of the form
\begin{equation}\label{eq:intro-signed}
 \int w(2M-D)
 \leq d\exp\!\left[-F_Q\!\left(\log\frac1d\right)\right]
\end{equation}
for nonnegative Lipschitz weights with fixed bounds. The explicit function
$F_Q$ grows faster than every fixed power of $\log T$ as $T\to\infty$.
\Cref{sx:signed} gives the precise statement, including the curvature error
for a reference surface of sufficiently high finite regularity. Retaining the sign permits a substantially
smaller error than an absolute-value estimate.

The proof of \eqref{eq:intro-signed} is an induction on $Q$. A
heat-potential construction turns a positive violation into a positive
measure on selected centers and scales. At each selected point it gives
an initial sign and a nonnegative cumulative integral along larger scales;
the selection includes vertical sides of the stopping boundary. For
configurations whose normal heights separate into groups, a stopping
argument reduces the estimate to stationary restrictions with smaller
integer counts. The resulting comparison controls tilt away from the groups.

Along a selected path, we take a supremum of damped signed-mass averages
over larger scales and compare it with a clipped squared-distance moment
from an unordered list of at most $Q$ representative normal heights.
A weighted combination of this moment and the cumulative contact integral
stays below the supremum even when representative heights merge. On an
interval without a merger, this comparison forces a quantitative amount
of tilt at many scales. Normal stationarity turns that tilt into a flux
pairing height displacement with radial tilt. A higher height moment
prevents all of that flux from concentrating on sets that have small base
volume but large local excess density. At each selected scale this yields
a base cube in which a quantitative fraction of the $\tau$-mass remains
below a density threshold at all smaller dyadic scales. Packing these
nested cubes shows that the union of their full innermost members has small
$\tau$-mass. A final charge bounds the contact measure by the mass of that
union, contradicting the original violation.

For regularity, fix a point where dilated varifolds converge to a
multiplicity-$Q$ plane
and the mass of other multiplicities has relative density zero. At each
nearby base point, we fit minimal graphs at successive scales, stopping
when a prescribed bound on the support height fails. The correction of
each fit is closely related
to the Jacobi regularization and nonlinear correction of
Brena--De Lellis--Franceschini \cite[Sections~2 and~6]{BDLF2025}.
A centered first-moment estimate, combined with elliptic regularity,
makes neighboring fits agree to second order in their height error.
The signed estimate is first used in a contradiction argument for a
sequence of failed fits normalized by their height errors. Together with
the limiting first-variation identities, its one-sided
inequality shows that the limiting squared-height measure is nonzero in
every interior subball. Strong control of the centered first moment,
together with the integer count $Q$, then forces that second moment to
come from separated projection entries of lower multiplicity on a definite
portion of each such ball. The regions
influenced by failures therefore have vanishing relative volume.

The compatible fits can now be glued into one fixed center of sufficiently
high finite regularity. The weighted integral of the squared mean curvature
and the square of its first derivative is bounded by the corresponding
squared-height integral. The support height is zero outside the regions
influenced by failures. Almgren developed the center-manifold approach
to higher-multiplicity regularity \cite{Almgren2000};
De Lellis--Spadaro construct their center by gluing smoothed local averages
\cite{DLS2016II}. For the inward step we use the comparison of height and
energy in the frequency method of this program
\cite[Section~3.4]{DLS2011}, including its cutoff form on a center manifold
\cite[Section~3]{DLS2016III}. The minimizing-current and
Dirichlet-minimizing hypotheses in the Almgren--De Lellis--Spadaro
regularity results are unavailable here, so the needed inequalities
are derived from stationarity and the signed estimate.

This second use of signed excess takes place at finite radii relative to
the fixed center. Interior balls retain the leading signed comparison
between projected mass and tilt. A separate radial test controls the thin
outer collar by a small multiple of the tilt in the radial direction.
It thereby preserves the leading tilt coefficient when the comparison is
inserted into the frequency identities. The resulting inward comparison
supplies height doubling without assuming it initially. Along a sequence
of radii with controlled frequency, the local mass estimate transfers the
vanishing relative volume of the failure regions to vanishing relative
projected mass. A higher height moment then
excludes nonzero height carried by those regions. The support lies on
the center in a neighborhood, and
stationarity gives constant multiplicity and makes the center a smooth
minimal graph.

\paragraph{Organization.}
\Cref{sec:local} establishes the first-variation, integer-count, and height
estimates. \Cref{sec:signed-statement} states the signed-excess theorem and
fixes its parameter hierarchy. \Cref{sec:contact} constructs the heat
contacts and reduces separated configurations to lower multiplicities;
\Cref{sec:excess} develops the outward comparison and controls mergers;
\Cref{sec:packing} closes the integer induction by packing witness scales.
\Cref{sec:centering} constructs the center from minimal fits and
multiplicity defects. \Cref{sec:inward} derives inward doubling,
forces the height to vanish, and proves the Hausdorff-measure conclusion
on the entire support. \Cref{sec:spheres} proves the round-sphere
consequence by the separate cone argument.

\subsection{Conventions}

All estimates are interior and local. Constants may depend on $m,n$, an upper
bound for the particular integer count $Q$, and the fixed coordinate geometry.
They are independent of the small scale and excess unless stated otherwise.
Fixed changes of concentric radii are allowed when positive interior room is
retained. A dot denotes differentiation with respect to the logarithm of the
radius. Every projection count includes the integer multiplicity.
We distinguish a positive tensor from the measure obtained by multiplying it
by $\mu$; projections of the latter to the base are stated when introduced.
\Cref{sec:signed-statement} specifies the quantitative hierarchy used in the
induction.
